\section{Residual conormals and exact local equations}\label{sec:residual-calculus}
We work over $\C$. Put $E_\mu=\bigoplus_\nu\Sym^2U_\nu$,
$M=Lc(k-c)$, $e=L\binom{c+1}{2}$, $p=2k-1$, and
\[
 a=e-p+1,\qquad b=Lc-3.
\]
When $e\ge p$, write $\mathscr D_L$ for the scheme of maximal minors of
$\mu:E_\mu\longrightarrow V_{2n}$, and $\mathfrak I$ for the incidence
$\ell\mu=0$ in $\Pj(V_{2n}^*)\times X_L$.
The following result is not restricted to a general point or to a stable
range of dimensions.

\begin{theorem}[Residual tangent sequence]\label{thm:residual-tangent}
Let $([\ell],U)\in\mathfrak I$, let $E=\operatorname{rad}H_\ell$, and put
\[
 W_\nu=U_\nu\cap E,\quad d_\nu=\dim W_\nu,\quad
 R_W=\sum_\nu W_\nu^2\subset V_{2n},\quad
 q=\binom{c+1}{2}.
\]
For each $\nu$ define
\[
 B_\nu:\operatorname{Hom}(U_\nu,V_n/U_\nu)\longrightarrow\Sym^2U_\nu^*,
 \qquad B_\nu(A)(u,v)=H_\ell(Au,v)+H_\ell(u,Av).
\]
Then $\operatorname{im}B_\nu$ consists exactly of the symmetric forms
whose restriction to $W_\nu\times W_\nu$ is zero. There is an exact sequence
\begin{equation}\label{eq:residual-tangent-sequence}
 0\longrightarrow\bigoplus_\nu\ker B_\nu
 \longrightarrow T_{([\ell],U)}\mathfrak I
 \longrightarrow R_W^\perp/\C\ell\longrightarrow0.
\end{equation}
In particular
\begin{equation}\label{eq:residual-tangent-dimension}
 \dim T\mathfrak I
 =M+p-1-\sum_\nu\left(q-\binom{d_\nu+1}{2}\right)-\dim R_W.
\end{equation}
If $\mu_U$ has corank one, these are also the tangent space and tangent
dimension of $\mathscr D_L$ at $U$; its conormal rank in $X_L$ is
\begin{equation}\label{eq:residual-conormal-rank}
 \sum_\nu\left(q-\binom{d_\nu+1}{2}\right)-p+1+\dim R_W.
\end{equation}
\end{theorem}
\begin{proof}
The map $V_n/U_\nu\longrightarrow U_\nu^*$ induced by $H_\ell$ has
image $(U_\nu/W_\nu)^*$. Indeed its transpose has kernel $W_\nu$;
quotienting by $U_\nu$ is permitted because $U_\nu$ is isotropic.
Choose a complement $Z_\nu$ of $W_\nu$ in $U_\nu$ and lifts of a dual
basis of $Z_\nu^*$. The coordinates of $A|_{W_\nu}$ in these lifts
prescribe every $W_\nu\times Z_\nu$ entry. The coordinates of
$A|_{Z_\nu}$ prescribe every symmetric $Z_\nu\times Z_\nu$ entry,
by taking one half of the desired symmetric matrix. This proves the
assertion about $B_\nu$, including its rank.

Differentiating $\ell(uv)=0$ gives
\[
 \dot\ell(uv)+B_\nu(\dot U_\nu)(u,v)=0.
\]
It is solvable for the $\dot U_\nu$ precisely when $\dot\ell$ annihilates
all $W_\nu^2$. Solutions for different contacts are independent. Taking
$\dot\ell$ modulo $\C\ell$ proves the exact sequence and its dimension.
If $\mu_U$ has corank one, invert a $(p-1)$-minor. Row and column
operations leave an identity block and one row of $a$ entries. The
normalized left kernel vector is uniquely solved by the identity block;
the remaining incidence equations are exactly these $a$ entries. Hence
the incidence projects isomorphically, as a scheme, to $\mathscr D_L$
on this chart. Subtracting the tangent dimension from $M$ gives the
conormal rank.
\end{proof}

\begin{proposition}[Local presentation at every corank]\label{prop:all-corank-schur}
At a point where $\mu$ has corank $\rho\ge1$, invert a
$(p-\rho)$-minor and write its matrix in blocks as
$\left(\begin{smallmatrix}A&B\\ C&D\end{smallmatrix}\right)$.
On this open chart, and also after completion of its local ring,
\begin{equation}\label{eq:all-corank-schur}
 I_p(\mu)=I_\rho(D-CA^{-1}B).
\end{equation}
The matrix on the right has size $\rho\times(e-p+\rho)$ and vanishes
at the given point.
\end{proposition}
\begin{proof}
Multiply by the invertible block matrices eliminating $C$ and $B$, and
then by $A^{-1}$ on the first block. The resulting matrix is
$\operatorname{diag}(I_{p-\rho},D-CA^{-1}B)$, in rectangular block
notation. Its maximal minors generate the indicated ideal. Invertible
row and column operations preserve determinantal ideals, and localization
and completion preserve this identity.
\end{proof}
This is a scheme presentation, not a replacement of its tangent cone by
minors of the linear part. Initial ideals can contain additional initial
forms obtained from cancellations. Nor does the presentation by itself
classify the associated primes of every excess-dimensional germ.

\subsection{The rank-two residual morphism}
Let $[\ell]$ have rank two and two distinct support points $x,y$.
Choose the two isotropic hyperplanes $H_+,H_-$ of $H_\ell$ and their
intersection $E=gV_{n-2}$, where $g$ vanishes at $x,y$.
Fix a sign pattern $\varepsilon$ and consider the open incidence on which
$U_\nu\subset H_{\varepsilon_\nu}$ and $U_\nu\not\subset E$.
\begin{proposition}[Dominance with affine fibres]\label{prop:residual-dominance}
For fixed $[\ell]$ and choices of its hyperplanes, the morphism
\begin{equation}\label{eq:residual-morphism}
 (U_\nu)_\nu\longmapsto
 \bigl(g^{-1}(U_\nu\cap E)\bigr)_\nu
 \quad\hbox{to }\Gr(c-1,V_{n-2})^L
\end{equation}
is smooth and surjective. It is locally an affine-space bundle of
relative dimension $L(k-c-1)$.
The same assertion holds relatively as the distinct support points and
their two hyperplanes vary, after trivializing the indicated bundles.
\end{proposition}
\begin{proof}
For a prescribed $W_\nu=gK_\nu$, an admissible $U_\nu$ is the inverse
image of a line in $H_{\varepsilon_\nu}/W_\nu$ projecting isomorphically
to $H_{\varepsilon_\nu}/E$. Its choices form a torsor under
\[
 \operatorname{Hom}(H_{\varepsilon_\nu}/E,E/W_\nu),
\]
an affine space of dimension $k-c-1$. Local choices of a splitting
trivialize this torsor. Intersections with $E$ have constant dimension
$c-1$, so the assignment is a morphism. Products give the result.
\end{proof}

\begin{corollary}[Exact singularity recursion on a signed stratum]
\label{cor:residual-singularity}
At a rank-two point as above, suppose that $\mu_U$ has corank one.
Let $\mu_K'$ be multiplication from $\bigoplus_\nu\Sym^2K_\nu$ to
$V_{2n-4}$ and put $\delta=\dim\operatorname{coker}\mu_K'$.
Then
\begin{equation}\label{eq:rank-two-tangent-excess}
 \dim T_U\mathscr D_L=M-b+\delta.
\end{equation}
If $4\le c<k$, $e\ge p$ and $b\le a$, this point is smooth in
$\mathscr D_L$ if and only if $\mu_K'$ is onto. In this range the
singular set on the stated corank-one open stratum is exactly the inverse
image of the residual multiplication failure locus.
\end{corollary}
\begin{proof}
Here $d_\nu=c-1$ and
$R_W=g^2\operatorname{im}\mu_K'$, of dimension $p-4-\delta$.
Equation~\eqref{eq:residual-tangent-dimension} gives the displayed formula.
The signed incidence has image of dimension $M-b$: a general first
$c$-plane in its secant hyperplane determines that hyperplane uniquely,
since the locus also lying in a second secant hyperplane has dimension
at most $3+c(k-2-c)<c(k-1-c)$. On the corank-one open set the local
signed incidence is embedded by the incidence isomorphism. Thus its
local dimension is $M-b$. Theorem~\ref{thm:global-intro} bounds the
dimension of the entire failure locus by $M-b$ when $b\le a$.
Consequently its local dimension here is $M-b$, and equality of tangent
and local dimensions is equivalent to $\delta=0$.
\end{proof}

\subsection{Projection and commutative algebra}
\begin{proposition}[From rank incidences to components]\label{prop:incidence-components}
Stratify the annihilator incidence by the rank $r$ of $H_\ell$ and by
$s_\nu=\dim(U_\nu/(U_\nu\cap\operatorname{rad}H_\ell))$.
Every irreducible component $C$ of the support of $\mathscr D_L$ is
dominated by the closure of an irreducible piece of one of these finitely
many strata. In particular a stratum of dimension less than $\dim C$
cannot dominate $C$.
If a dominating incidence has dimension $\dim C$, and its generic fibre
contains a nonempty open subset of the projective annihilator space,
then that annihilator space is a point and its projection is birational
in characteristic zero.
\end{proposition}
\begin{proof}
The projective incidence maps onto the failure locus. Intersect the
images of its finitely many locally closed strata and their irreducible
pieces with $C$. These constructible sets cover $C$; at least one is
dense. The dimension theorem gives the first assertion and implies that
the generic fibre of an equal-dimensional dominating piece is
zero-dimensional. The projective annihilator space is the projectivization
of a vector space, not an arbitrary finite collection. A nonempty open
subset of a positive-dimensional projective space has positive dimension.
Thus the open-fibre hypothesis forces the projective space to be a point.
On a corank-one chart the incidence is the graph of its uniquely solved
annihilator. This proves birationality, including the scheme-theoretic
generic-point assertion.
\end{proof}

\begin{proposition}[Expected grade and absence of embedded primes]
\label{prop:expected-grade}
Suppose $4\le c<k$, $e\ge p$, and $a\le b$. On every nonempty affine
Grassmannian chart meeting $\mathscr D_L$, its maximal-minor ideal has
height and grade $a$. The Eagon--Northcott complex resolves its quotient,
which is Cohen--Macaulay of pure codimension $a$ and has no embedded
associated primes.
\end{proposition}
\begin{proof}
The chart ring is regular. The global codimension theorem gives height
at least $a$, whereas the maximal-minor height bound gives height at most
$e-p+1=a$. In a regular ring grade equals height. The Eagon--Northcott
grade criterion gives a free resolution of length $a$ and the
Cohen--Macaulay conclusion \cite{EN,Kustin}. Localization gives the same
conclusion at every prime over the ideal. A Cohen--Macaulay quotient of
this expected dimension is unmixed and has no embedded associated
primes. No independence assumption on the matrix entries is used.
\end{proof}
\begin{proposition}[The reducedness step]\label{prop:generic-reducedness}
In the situation of Proposition~\ref{prop:expected-grade}, if
$\mathscr D_L$ is generically reduced on each of its irreducible
components, then it is reduced. If its support is also irreducible, it
is integral.
\end{proposition}
\begin{proof}
A Noetherian ring without embedded associated primes satisfies $S_1$.
Generic reducedness is $R_0$. The criterion $R_0+S_1$ for reducedness
therefore applies \cite[Tag 031O]{Stacks}.
Alternatively a nonzero nilradical would have an associated prime;
without embedded primes this would be a minimal prime, contradicting
generic reducedness. The last assertion is the definition of an
integral scheme.
\end{proof}
